\section{Организация и проведение похода}
\subsection{Цели и задачи маршрута. Выбор нитки маршрута}
При разработке и планировании маршрута я, как руководитель, руководствовался следующими соображениями:
\begin{enumerate} 
	\item \textbf{Высотность, транспортная доступность горного района и стоимость трансфера.}
	Подавляющая часть группы (пять человек из девяти) имела официальный или неофициальный опыт горных походов до 2 к.с по Кавказу \cite{Snegovskaya2024} или Алтаю, и вся группа имела высотный опыт ночёвки не менее 2900 м \cite{ostapiv2025}. Это позволило выбрать в качестве района новичкового похода не Гвандру или Архыз, как это обычно бывает, а новый для участников и руководителя горный район Центральной Азии~--- Тескей-Ала-Тоо. Высоты в этом районе, в среднем, на 500 м больше, чем на Кавказе, а его транспортную доступность можно назвать приемлемой для новичковых походов: 4 часа самолётом до Бишкека и еще 7~--- до места старта в горах. Стоимость трансфера выше, чем на Кавказе, но ниже, чем на Алтае. 
		
	\item \textbf{Концентрированность препятствий.}
	Дополнительным фактором в сторону выбора Тескей-Ала-Тоо послужило также и то, что в отличие, например, от Алтая, хребтовка района (основной хребет и серия параллельныз отрогов), как и на Кавказе, позволяет осуществить подход под многие перевалы в течение одного дня, миновав долгие и занудные забеги по долинам, и, как следствие, поддерживать интерес группы на приемлемом уровне.
	
	\item \textbf{Разнообразие рельефа.} 
	С методической точки зрения, а также, опять-таки, для поддержания интереса группы, хотелось продемонстрировать участникам как можно больше разнообразных типов рельефа. Район позволяет в походах 1 к.с. продемонстрировать все виды рельефа, за исключением, пожалуй, снежного. <<Изюминкой>> района является возможность пересечения главного хребта через перевалы 1А (разительно отличающихся от классических <<единичек А>>) с осмотром высокогорных болот~--- сыртов.
	
	\item \textbf{Эффект кульминации.}
	У меня как у руководителя было глубокое убеждение, что первый поход должен обладать понятным, с позволения сказать, сюжетом и иметь свою кульминацию. В нашем случае таким сюжетом было движение через отроги главного хребта с преодолением разнообразных перевалов (последовательно: травянисто-осыпной, ледовый, травянистый), а затем~--- кольцо через главный хребет с <<неклассическими>> перевалами на южную сторону, восхождением на обзорную точку~--- вершину Марс. Конец маршрута~--- забег по среднегорью с финишем на локальной достопримечательности <<Красные скалы>>~--- также должен был добавить в поход разнообразия и красоты.
	
	\item\textbf{Помощь в планировании}
	Фактор, который формально не был в списке определяющих критериев, но по факту являлся таковым~--- это искренняя заинтересованность и помощь в планировании и организации маршрута от человека, через которого организовывалась логистика похода~--- Юрия Траченко. За эту помощь выражаю ему огромную благодарность.
	
\end{enumerate} 

\subsection{Логистика}
Подъезд осуществлён на поезде 033М Москва--Владикавказ до станции Минеральные Воды (прибытие в 03:40). Стоимость проезда на август 2024 г. составляла 7800~\faRub, купе (обратно – 4700~\faRub, плацкарт). От Минеральных Вод до аула Верхний Учкулан (время в пути 4 часа) добирались на трансфере, заказанном через Саракуева Бориса (89289503868, 89298843175,  \href{mailto: bezonec@list.ru}{bezonec@list.ru}). Стоимость трансфера трансфера туда составила 18000~\faRub, обратно (от поляны Азау) — 15000~\faRub. Стоимости доставки забросок в т/б Глобус и а/л Узункол составили 4000~\faRub~и 6000~\faRub.
Коллективный пропуск в пограничную зону КЧР был оформлен за 4 месяца до начала похода через электронную почту пограничного управления ФСБ по КЧР~--- \href{mailto: pu.kcherkes@fsb.ru}{pu.kcherkes@fsb.ru} и отправлен письмом по указанному адресу. Пропуск в КБР не требуется, так как пер. Хотютау в 2023 году был исключён из пограничной зоны \cite{order_kbr}.

\subsection{Аварийные выходы из маршрута и его запасные варианты}

\textbf{Аварийными выходами} с маршрута являлись:
\begin{itemize}
%	\item \textbf{5 августа:} д.р. Кичи-Кызыл-Суу – с. Кызыл-Суу;
%	\item \textbf{6 августа:} д.р. Джукучак – д.р. Джууку – с. Саруу;
%	\item \textbf{7, 12, 13 августа:}д.р. Ашу-Кашка-Суу – д.р. Джууку – с. Саруу;
%	\item \textbf{8 августа:} д.р. Иттиш – д.р. Джууку – с. Саруу
	\item \textbf{03–05.08:} Спуск по д.р. Чон-Кызыл-Суу к п. Кызыл-Суу;
	\item \textbf{05–06.08:} Спуск по д.р. Кичи-Кызыл-Суу к п. Кызыл-Суу;
	\item \textbf{06–07.08:} Спуск по д.р. Джуукучак к слиянию с р.Джууку;
	\item \textbf{07–09.08:}Спуск по д.р. Ашукашкасу и р. Джууку к слиянию с р. Джуукучак;
	\item \textbf{09–11.08:}Спуск по д.р. Ит-Тиши(юж.) в д.р. Ара-Бель-Суу к дороге на Кумторский рудник.
	
	
\end{itemize}
\textbf{Запасными вариантами} маршрута являлись:

\begin{itemize}
	\item Отказ от восхождения на в. Марс

\end{itemize}


\subsection{Изменение маршрута и их причины}
Маршрут пройден без изменений.
\subsection{Обеспечение безопасности на маршруте}
Группа была зарегистрирована в отделении МЧС Кыргызстана по Иссык-Кульской области (телефон +996555004214, связь по WhatsApp). За 10 дней до похода были переданы сведения о составе группы, сроках и маршруте, в ответ был получен регистрационный номер и просьба проинформировать дежурного о начале и завершении похода.

Для регулярного обмена звонками и сообщениями были арендованы спутниковые телефоны системы Iridium: модель Iridium 9575  для нашей группы и Iridium 9555 для дружественной группы Алексея  Остапива. Стоимость аренды одного телефона в компании ООО <<TECC.KOM>> (www.tecckom.ru) на 18 дней составила 17100~\faRub, активация СИМ-карты ~--- 4800~\faRub. По окончании аренды эти суммы вместе с услугами связи (стоимость исходящей минуты 130~\faRub, стоимость СМС ~--- 80~\faRub) вычитаются из предварительно внесенного залога 50000~\faRub. 

В целом, от использования спутниковых телефонов остались смешанные ощущения. Предварительное тестирование аппаратов в Москве показало, что звонки со спутников на обычные телефоны проходят через раз (это же было подтверждено в походе, когда в условиях, свободных от глушения спутниковой связи, вообще не получилось дозвониться до одного из координаторов); при этом связь спутниковых телефонов друг с другом была стабильной. Исходящие СМС-сообщения отправлялись, но не доходили до адресатов, с входящими проблем не возникало. Подобную ситуацию можно описать словами <<не баг, а фича>>, поскольку данные особенности могут быть связаны с настройками серверов Iridium или телефонов адресатов, причем заранее предугадать поведение в той или иной ситуации невозможно. Поэтому, когда в походе нет необходимости регулярно быть на связи с другой группой, рекомендуется использовать спутниковый треккер.
%\alert{Дима, это с тебя Для регулярного обмена сообщениями, отслеживания положения группы на карте, а также возможности экстренной связи, в группе имелся спутниковый треккер IRIDIUM Rockstar 360. Стоимость аренды треккера в <<Альпиндустрии>> на 15--21 день составила 7100~\faRub, залог~--- 50000~\faRub. Нам повезло попасть на демострационный период тарифа треккера, в связи с чем все сообщения были для нас безлимитны и бесплатны. Предварительное тестирование треккера в Москве показало, что спутниковые сигналы в столице эффективно глушатся: сообщения приходили не чаще раза в сутки. В походе с приёмом и отправкой сообщений и координат на сервер проблем не возникало, среднее время отправки составляло 30 минут.}

Каждый участник самостоятельно оформлял на себя индивилуальный страховой полис. Выбрали страховую фирму <<Согласие>>, ассист Balt Assistance, опция <<Треккинг свыше 1500 м>>, размер страховой защиты 35000~USD,  вид отдыха <<Спорт Экстрим>>. Стоимость полиса составила 8234~\faRuble~с человека.


\subsection{Перечень наиболее интересных природных и исторических объектов, занятий на маршруте}
\begin{enumerate}[noitemsep,topsep=0pt,parsep=0pt,partopsep=0pt]
	
	\item Широкие красивые долины северной стороны хребта: Чон-Кызыл-Суу, Киче-Кызыл-Суу, Ашукашкасу, Джукучак;
	\item Сырты на южной стороне хребта;
	\item Цепочка озёр Кашкасу с чистой водой;
	\item Вершина Марс, с которой открываются виды на Акширак, Кумтор, сырты и, в хорошую погода, на семитысячники;
	\item Перевал Кашкасу, через который в былые годы перегоняли лошадей. Следы этого остаются в ущелье до сих пор, жутковато и атмосферно;
	\item <<Красные скалы>> на слиянии р. Джууку и Джукучак;
	\item Иссык-Куль.
\end{enumerate}

\paragraph{Темы практических занятий:}

\begin{itemize}
	\item Техника передвижения по травянисто-осыпным склонам;
	\item Техника несложных бродов поодиночке;
	\item Техника установки лагеря в сильную непогоду.
	
\end{itemize}

\newpage
\subsection{Развёрнутый график движения}
\begin{table}[h!]
	\centering
	\resizebox{0.95\textwidth}{!}{%
		\begin{tabular}{|>{\centering\arraybackslash}m{0.045\linewidth}
				|>{\centering\arraybackslash}m{0.02\linewidth}
				|>{\centering\arraybackslash}m{0.43\linewidth}
				|>{\centering\arraybackslash}m{0.09\linewidth}
				|>{\centering\arraybackslash}m{0.1\linewidth}
				|>{\centering\arraybackslash}m{0.05\linewidth}
				|>{\centering\arraybackslash}m{0.09\linewidth}
				|>{\centering\arraybackslash}m{0.13\linewidth}|}
			\hline						
			Дата	&	\begin{turn}{90}День\end{turn}	&	Участок маршрута	&	Км с $k=1.2$	&	Набор /сброс, м	&	ЧХВ	&	Высота ночёвки, м	&	Способы передвижения	\\
			\hline
			
			03.08	&	1	&	м.н.~--- кур. Джилысу ~--- ФГС (д.р. Саватор)&	6.0	&	$+550$\newline$-260$	& 5:07	&	2623	& Пешком	\\
			\hline
			04.08	&	2	&	ФГС ~--- д.р. Чон-Кызылсу ~--- д.р. Саватор ~--- д.р. Каратакия ~--- м.н. 	&	4.8	& $+788$\newline$-0$		& 3:53		& 3430		&	Пешком	\\
			\hline
			05.08	&	3	&	м.н. ~--- д.р. Каратакия ~--- ~--- \textbf{пер. Каратакия (1А, 3800)}~--- - д.р. Кичи-Кызылсу – м.н. м.н.	&	5.8	& $+439$\newline$-531$		& 4:37	& 3280		&	Пешком	\\
			\hline
			06.08	&	4	& м.н.  ~---\textbf{пер. Кашкатор Сев. (3929, 1А)}~---  д.р. Джуукучак ~--- м.н.	&	12.7	&$+660$\newline$-1005$		& 7:47		& 2900		&	Пешком	\\
			\hline
			07.08	&	5	& м.н.  ~---\textbf{пер. Ашутор Западный (3700, 1А)}~---  слияние р. Ашукашкасу и р. Джууку ~--- м.н., заброска	&	11.4	& $+750$\newline$-1183$		& 5:53	& 2500		&	Пешком	\\
			\hline
			08.08	&	6	&	м.н. ~--- д.р. Джуукучак ~--- д.р. Иттиши ~--- м.н.	&	8.4 	& $+809$\newline$-0$		& 4:49		& 3285		&	Пешком	\\
			\hline
			09.08	&	7	&	м.н.  ~---\textbf{пер. Иттиш (3895, 1А)}~---   м.н.	&	11.8	& $+569$\newline$-140$		& 6:23		& 3870		&	Пешком	\\
			\hline
			10.08	&	8	&	\textit{\textbf{Днёвка}}	&	0	& $+0$\newline$-0$		& 0		& 3870		&	Пешком	\\
			\hline
			11.08	&	9	&	м.н. ~--- д.р. Ит-тиши ~--- Сырты ~--- оз. Кашкасу ~--- \textbf{в. Марс (н/к, 4345,	рад.)} ~--- м.н.	&	25.8	& $+669$\newline$-522$		& 9:08		& 3900		&	Пешком	\\
			\hline
			12.08	&	10	&	м.н.  ~--- \textbf{пер. Кашкасу (3890, 1А)}~---  д.р. Ашуу-Кашка-Суу ~--- м.н.	&	11.2	& $+0$\newline$-900$		& 5:43		& 3000		&	Пешком	\\
			\hline
			13.08	&	11	&	м.н. ~--- слияние р. Ашуу-Кашка-Суу и р. Джууку ~--- д.р. Джууку – слияние р. Джуукучак и Джууку ~--- (финиш) &	25.8	& $+0$\newline$-1010$		& 5:00		& 1990		&	Пешком	\\
			
			\hline
			\multicolumn{3}{|c|}{\textbf{\textit{\Large{Итого:}}}} & \large{\textbf{123.6}} & \large{$\mathbf{+5234}$\newline$\mathbf{-5551}$	}	& \multicolumn{3}{c|}{\large{\textbf{58:20}\newline\textbf{2д 10ч 7мин}}} \\
			\hline
		\end{tabular}
}	
	
\end{table}



\clearpage